% Generated draft: requires mathematical and source review.
\documentclass[11pt]{article}
\usepackage[margin=25mm]{geometry}
\usepackage{fontspec}
\setmainfont{Latin Modern Roman}
\usepackage{amsmath,amssymb,mathtools,unicode-math}
\setmathfont{Latin Modern Math}
\usepackage{enumitem,needspace,xurl,hyperref}
\setlength{\parindent}{0pt}
\setlength{\parskip}{5pt}
\setlength{\emergencystretch}{3em}
\newcommand{\N}{\mathbb{N}}
\newcommand{\Z}{\mathbb{Z}}
\newcommand{\Q}{\mathbb{Q}}
\newcommand{\R}{\mathbb{R}}
\newcommand{\C}{\mathbb{C}}
\newcommand{\F}{\mathbb{F}}
\newcommand{\A}{\mathbb{A}}
\newcommand{\PP}{\mathbb{P}}
\newcommand{\E}{\mathbb{E}}
\newcommand{\eps}{\varepsilon}
\newcommand{\dd}{\,\mathrm{d}}
\begin{document}
\section*{UnsolvedMath ID: 30002873}
\textbf{Source problem code:} catalogue entry
\subsection{Definitions and mathematical statement}
\textbf{Source question.} Let \(\Delta\subset\Z^n\) be a Newton polyhedron and let \(H\) be a generic hypersurface with Newton polyhedron \(\Delta\). Given a dominant rational map
\[
f:H\dashrightarrow\PP^{n-2},
\]
determine the genus of its generic fiber. The target requires an integer \(n\geq 2\), although this restriction is not explicit in the source. The map \(f\) is part of the given data; the record does not specify a particular map or assert that the genus depends only on \(\Delta\).

\textbf{Newton data.} For a Laurent polynomial \(P=\sum_{a\in A}c_a t_1^{a_1}\cdots t_n^{a_n}\), with finite support \(A\subset\Z^n\) and nonzero coefficients \(c_a\), its Newton polytope is \(\operatorname{conv}(A)\subset\R^n\). The source instead calls \(\Delta\subset\Z^n\) a Newton polyhedron. It does not clarify whether this denotes the support, a lattice polytope, or another polyhedral convention. It also leaves the ground field, ambient variety, coefficient family, and meaning of generic unspecified.

\textbf{Conditional interpretation of genus.} If the intended setting is an integral hypersurface \(H\) of dimension \(n-1\) over a field \(k\), set
\[
K=k(\PP_k^{n-2}),\qquad \eta=\operatorname{Spec}K.
\]
On a dense open subset \(U\subset H\) where \(f\) is a morphism, its generic fiber is \(U_\eta=U\times_{\PP_k^{n-2}}\eta\), a curve over \(K\). If this curve is geometrically integral, one possible intended invariant is the geometric genus
\[
g(H,f)=\dim_{\overline K}H^1(\widetilde C,\mathcal O_{\widetilde C}),
\]
where \(\widetilde C\) is the smooth projective model of \(U_\eta\times_K\overline K\). Under this interpretation, the task is to determine \(g(H,f)\) from the specified hypersurface and rational map.

These are conditional definitions, not additional hypotheses supplied by the source. Dominance alone does not ensure geometric integrality of the generic fiber. The record does not specify geometric versus arithmetic genus, a compactification, or a convention for geometrically reducible or nonreduced fibers. Its notation \(f^{-1}(x)\), with \(x\in\PP^{n-2}\), also leaves unresolved whether it means the fiber over the generic point or a fiber over a sufficiently general geometric point. Consequently, the supplied question requires these qualifications before it has a unique precise formulation.
\subsection{Short English statement}
Given a generic hypersurface with prescribed Newton data and a dominant rational map to \(\PP^{n-2}\), determine the genus of the generic curve fiber. The field, Newton convention, meaning of generic, and intended genus require clarification. The supplied record labels the problem open; this draft does not independently verify that status.
\subsection{Sources}
\begin{itemize}
\item Local registry: \texttt{data/unsolvedmath/problems.json}, permanent UnsolvedMath ID \(30002873\), source code \texttt{catalogue entry}, title \emph{Fiber Genus of Rational Maps from Newton Hypersurfaces}; collection: Oberwolfach Reports Open Problems. No individual proposer is supplied.
\item Supplied report attribution: \emph{Tropical Aspects in Geometry, Topology and Physics (2016)}, Oberwolfach Reports, pp.\ 1255--1316, DOI \texttt{10.4171/owr/2015/23}. \url{https://doi.org/10.4171/owr/2015/23}. The background describes Oberwolfach Reports as records of mathematical workshops and their open-problem discussions. The record gives proposed year 2015, while its report citation includes 2016; this discrepancy is unresolved.
\item Supplied literature source: Oberwolfach Report 23. \url{https://doi.org/10.4171/owr/2015/23}.
\item Supplied literature source: Stapleton--Ullery, \emph{Degree of Irrationality of Hypersurfaces}. \url{https://arxiv.org/abs/1803.02957}.
\item Literature assessment dated exactly \texttt{2026-08-21}: supplied dated curation, not independently verified. It labels the question open and states that degree-of-irrationality results constrain related maps and fibers but do not provide the requested general genus formula. The background characterizes this assessment as dated literature triage derived from supplied OWR curation and requests independent verification of time-sensitive status and cited claims.
\item The registry supplies statement status \texttt{corrected\_verified}, described as preserving mathematical content while repairing notation, typography, or omitted definitions for publication. This is a supplied editorial label, not independent mathematical verification by this draft.
\end{itemize}
\end{document}
